\section{Diverse tips}
\subsubsection{I er stadig på scenen selvom I ikke spiller}
%Husk, at publikum altid kan se og høre jer, når I er på bandscenen. UP1 er  meget lydt, så småsnak på bandscenen kan godt høres blandt publikum.
Husk at I også er på scenen under sketches!
Bandet opholder sig gennem hele revyen på bandscenen.
Det betyder i praksis at selvom der under en sketch primært er fokus på scenen vil der på alle tidspunkter sidde publikum på pladser i salen som kan se bandet.
Hold jer derfor i ro så fokus forbliver på den sketch jeres scenegængere er ved opføre.
Vær dertil opmærksomme på jeres generelle udstråling på scenen.
Selvom I har set sketchesne hele ugen og efterhånden er trætte af den samme joke der bliver fortalt igen er det vigtigt at holde en ordentlig mine som ikke er forstyrrende hvis/når publikum ind imellem får skelet til bandscenen.

Noddling på, og stemning af elguitarer kan også høres i salen, selvom man er mutet.

Specifikt til blæsere: Mikrofoner er ind imellem tændt kort tid inden en sang starter/ efter en sang slutter.
Mikrofonerne kan opfange snak i nærheden og det er derfor vigtigt ikke at sniksnakke med hinanden hvis det ikke er yderst nødvendigt.

-Christine Mehlsen

\subsubsection{Nodelamper tiltrækker opmærksomhed}
Når der er mørkt i salen, er nodelamper på bandscenen meget iøjefaldende for publikum, fordi lamperne peger ned imod dem. Det kan være en fed effekt, at der er mørkt på bandscenen indtil lyset fra teknikken tændes. Hvis man memoriserer de første par takter af hvert nummer, kan man vedligeholde en skarp overgang fra mørke til lys når bandet begynder at spille.

\subsubsection{Intonation af akkorder}
Hvis sangere og/eller instrumenter der intonerer skal spille hver deres tone i en akkord, kan det være svært at få til at stemme. Et godt redskab til dette er at spille (langsomt) indtil dén akkord der lyder falsk, og så stoppe på den, og holde den. Lyt til den og korriger intonation. Prøv så det samme igen og se om i kan få den til at stemme lige når den slås an. Fortsæt derefter evt. indtil næste akkord og lav samme øvelse.

\subsubsection{Spil hele året}
Se, om der er projekter i løbet af året I kan samle bandet til, ud over revyen. Hvis jeres studie har en årsfest, kan i måske spille til denne som partyband?

\subsubsection{Tag ind og se andre revyer, og network med andre revybands}
Tag med ind og se de andre revyer i løbet af året! Det er sjovt og inspirerende. Lyt til hvordan de lyder i rummet, hvordan de arrangerer deres musik, og hvordan de arbejder sammen med scenen. Snak med dem om de ting de har lavet, og hvilke udfordringer de her. Ofte får man fede ideer til ting man selv kan gøre i sin egen revy. Af og til får man også gode ideer til ting man kan lade være med.

\subsubsection{Søg penge til grej}
Nogle gange kan man få penge til indkøb af inventar til dyrestalden. SciencePuljen har ofte givet her. Klaveret, barytonsaxen, PA'et og bassen/basforstærekeren er købt af fondspenge, så sørg for at søge!

\subsubsection{Materialeudvælgelse: kreativ vision}
Skal revybandet tjene intruktørenes kreative vision, eller skal bandet være med til at drive den kreative vision? Jeg hælder helt klart mere til det sidste. Advokér til materialeudvælgelsen for numre der er fede at spille, der passer godt til bandet, eller som du eller andre har en kreativ vision for.

\subsubsection{Materialeudvælgelse: hvad der fungerer som livemusik}
Den øvrige kreative ledelse af revyen mangler ofte fornemmelse for, hvad der fungerer godt som livemusik.
Jeg tænker, at det er på sin plads, at bandbossen til dels bærer den hat.
Ikke, at man skal gå at lege kong gulerod, men en vurdering af potentiale som livenummer er alt andet lige en fin tiebreaker, hvis det er svært at vælge imellem to numre.

-Ulrik "Larmcentral" Thinggaard Hansen

\subsubsection{Materialeudvælgelse: hvad er scenegængere interesserede i at opføre?}
Det materiale der scorer højt til sangudvælgelse eller vurderes højt af de kreative udvalg svarer ikke altid til, hvad scenegængere har lyst til at opføre.
Klassikeren er, at folk gerne vil stemme rapmusik ind, men de fleste, der står på scenen, vil gerne synge og vil ikke rappe.
Det burde naturligt sætte et loft over, hvor mange raps, der er i en revy.

-Ulrik "Træet" Thinggaard Hansen

\subsubsection{Bandet er (ikke) en kult}
Nogle gange kan et revyband virke som en kult: De er sære, taler i koder og holder sig mest for sig selv.
Det er selvfølgelig vigtigt at bandet er sammentømret og trygge med hinanden.
I midlertid burde man alligevel gøre en indsats for at integrere bandet med resten af revyen, for eksempel ved ikke altid at have et bandbord til aftensmad.
Det er også værd at være bevidst om at bandet er den eneste afdeling af revyen, der har adgangskrav og optagelsesprøve.
Det medfører en større risiko for at blive opfattet som elitær, hvis ikke man gør en indsats for det modsatte.

\subsubsection{Hjælp nye bandlemmer med at forstå hvad fællesbeskeder betyder for bandet:}
Ofte bliver fællesbeskeder givet med henblik på scenegængere.
Hjælp nye revyster i bandet med at få den nødvendige information, også  om de små ting.
Gør dem f.eks. opmærksom på jordbærdrinks/minirevy, mind dem om at skrive sig på aftensmadsarket i tide osv.

-Christine Mehlsen

\subsubsection{Tag kontrol når der øves på scenen}
De sidste par dage af revyforløbet øver man for det meste på scenen. Øvetid på scenen er både knap og efterspurgt, og når der både er skuespillere/sangere, instruktører, band og teknik involveret, er det mange personers tid der er på spil. Derfor er det vigtigt at være skarp for at få udnyttet tiden. Som bandboss har man en unik position til at facilitere dette: fra bandscenen kan man nemt se og kommunikere med scenen, instruktører, teknik og band (inklusiv trommeslager). Derudover har man som musiker ofte mere erfaring med ensemble-øvning end de andre afdelinger.
Udnyt disse styrker: Identificér hvilke specifikke dele af nummeret der behøver øvning og arbejd målrettet med disse. Undgå hele gennemspilninger hvis ikke det er nødvendigt. Hold kommentarer til det konstruktive og nødvendige og hold skruen i vandet: Hvis blot alle er enige om, hvad der skal ske, er det vigtigste at få flere gentagelser af de svære passager ind.
Vær ikke bange for at tage diktator-hatten på her. Det er i virkeligheden rart at der er nogen der tager styring i en ellers kaotisk situation, og folk vil takke dig hvis du tager styringen og giver øvningen retning.

\subsubsection{Tag kontrol over publikum til forestillingerne}
Det er et nemt publikum at få til at deltage. Brug dette på strategiske(musikalske) tidspunkter. Hvis man klapper med armene over hovedet og synligt henvender sig til publikum klapper de typisk med. At få publikum til at klappe på alle 4. eller 8. dele i et buildup kan virkelig pumpe energi ind i det.\\

\subsubsection{Brug stortrommen til at kontrollere publikums tilråb}
Ofte begynder publikum at råbe dumme slogans i takt. Det er for det meste sjovt, men nogle gange gør det en sketch for lang eller kommer i vejen for flowet. Hvis man bruger stortrommen til at trampe med på råbet kan man give det et skub de første 2-3 omgange. Hvis man derefter stopper, stopper publikum for det meste også instinktivt: når kicken pludselig falder ud, bliver energien i råberiet pludselig flad.

\subsubsection{Lyd i Store UP1}
Dette er en hel bog for sig, som jeg håber, nogen vil skrive. Nogle ting at være opmærksom på: baslyd skal formes, hvis den skal kunne være høj nok, uden at blive overvældende. En nem løsning er at køre signalet igennem en amp og sende direct out fra ampen til pulten. Blæsere klinger ofte i samme spektrum som sangerne.

\subsubsection{Effektiv oprydning efter revy}
Typisk er nedtagning af scenen det sidste der bliver færdigt under en revyoprydning. Derfor er det vigtigt af komme i gang så hurtigt som muligt, og noget af det første der skal ske, er at bandscenen skal tages ned. Bandets vigtigste opgave er derfor at rydde bandscenen så hurtigt som muligt.
Start med at tage instrumenter ned, og læg dem ud i dyrestalden. Fjern derefter andet gear, nodestativer, monitorsystemer, kabler, mikrofoner og deslige. Det er fint at lægge på hovedscenen, så længe det ikke kommer i vejen for teknikken.
Tag trommeburet ned når der er plads til det, og tag derefter trommesættet ud i dyrestalden.
Så er bandscenen klar til at blive taget ned.

\subsubsection{Om valg af nodestativer og montering på rælingen}
Der er som udgangspunkt to muligheder for valg af nodestativer:
\begin{enumerate}
    \item De store sorte fra dyrestalden:\\
        Pros: giver et ensformigt look fra publikum\\
        Cons: tager meget af udsynet og gør det sværere at se bandet.
    \item Mindre nodestatider som kan foldes sammen\\
        Pros: fylder mindre og tager dermed mindre udsyn fra publikum\\
        Cons: kan se rodede ud med mange løsblade på scenen.
\end{enumerate}

Begge typer monteres bedst med gaffatape på rælingen.
Hvis stativet har ben/fødder, så slå dem sammen inden opsætning, men lad være med at afmontere dem fra resten af nodestativet.
Erfaring har vist at de er svære at sætte sammen igen på bagkant.

- Christine Mehlsen

\subsubsection{Sådan hjælper du svage sangere (i praksis)}
Det hyppigste problem som sangere oplever er problemer med at finde ind på deres stemmer. Når først de er der, klarer de det som regel fint.
Har man en sanger der har svært ved at finde sin starttone i et bestemt stykke, kan man hjælpe ved at få et instrument til at spille den tone sangerens frase starter på, umiddelbart inden. På den måde får sangeren givet sin tone direkte, uden at skulle inferere den fra den harmoniske kontekst.
Det lyder måske som en klodset løsning, men det kan ofte gøres på en måde der lyder helt naturlig, for eksempel ved at lægge tonen øverst i sin guitar/keyboard voicing.
Hvis sangeren har meget svært ved det, kan man også lægge hele melodien i en af instrumenternes rolle.
Disse teknikker fungerer også fint til at hjælpe korsangere med at finde/holde deres stemme.

\subsubsection{Sådan hjælper du svage sangere (i princippet)}
Den største hjælp man kan give uerfarne sangere er, at hjælpe dem med at øve godt.
Nogle sangbosser vil mene, at det vigtigste for at en sanger kan præstere, er deres selvsikkerhed.
Selv da er god øvning det vigtigste redskab, da konkrete og reelle succesoplevelser er det vigtigste materiale til opbygning af selvsikkerhed.

\subsubsection{Når flere revyer øver samtidig}
Når flere revyer ligger tæt op ad hinanden, og har øveweekend samtidig, er dyrestalden ikke nok. Hør om nogen i bandet (eller revyen generelt) bor på et kollegie der har et øvelokale, I kan bruge. Kollegiegården, Egmont og Studentergården har alle øvelokaler, og ligger tæt på DIKU, så det også kan give mening for sangere at komme forbi.

\subsubsection{Hav en plan for indlæring af krævende materiale}
Udfordrende materiale kan nemt komme til at sluge for meget tid. Sørg for at have en plan for hvordan krævende materiale skal indlæres. Bryd det ned i mindre stykker, identificér de mest udfordrende og fokuser på dem individuelt. Opdel evt. bandet hvis forskellige dele er svære for forskellige sektioner, og øv parallelt. På den måde bruger man bruger mindre tid på at vente på hinanden.

\subsubsection{Når nogen forsvinder fra din monitor, når andre begynder at spille}
Monitor-controllerne vi har på bandscenen har en global kompressor/limiter. Hvis der er stor forskel i niveau mellem forskellige signaler og der ikke er nok headroom, kan kompressoren føre til at svagere signaler drukner når stærkere signaler kommer på. Prøv at slå den fra eller sænke ratio/hæve threshold. 

\subsubsection{Brandsikkerhed og bandscenen}
Bed scenemester om at sørge for, at alle rekvisiter der stilles bag bagtæp-
pet har en ansvarlig der fjerner denne i tilfælde af brand. Således kan
bandet også komme ned af bandscenen hvis det bliver nødvendigt.

\subsubsection{Husk at holde pauser}
De flest sangere/instruktører vil gerne have så meget tid med bandet som muligt. Det fører nemt til at bandet har øvesegmenter uafbrudt hele dagen. Sørg for at der bliver indlagt deciderede pausesegmenter til bandet. De 3 minuter det tager at skifte sanger-hold er ikke nok.